\chapter{总结与展望}
\label{ch:concl}

\section{工作总结}

本文以 Bender 等人提出的最优时间/空间权衡框架\cite{bender2022otsh} 为理论基础，实现了基于 k-kick 树的分层哈希表的结构与算法设计，并据此实现简化原型、开展性能实验。实现中的主要难点在于将 mini-array、k-kick 与 local query router 串联合并为统一系统，并在 facility--cubby--slot 分层组织下保持辅助结构同步不变量。在此基础上，本文进一步加入冷热分层键存储方案，用内存 tail cubby 存储近期数据，用 SQLite3 管理下沉后的冷数据\cite{gaffney2022sqlite,dong2021rocksdb}。

设计部分采用 facility--cubby--slot 分层组织。更新路径主要依赖 mini-array 维护变长元数据，并通过 k-kick 控制插入过程中的有限元素交换。查询路径由 local query router 定位候选槽位，再结合商化压缩信息恢复和校验键。每次发生搬移时，槽位数组、mini-array \(B\)、\(M\) 与 local query router 都需要同步更新。

实验测量分为两部分。第四章表~\ref{tab:exp-n-fixed}--\ref{tab:exp-tier-fixed} 及图~\ref{fig:exp-n-latency} 检查基础多层级结构的功能与维护成本；踢出深度变化时，单次搬移次数仍不超过对应的 \(k\)，与前文交换上界一致。第五章表~\ref{tab:hotcold-memory-latency}--\ref{tab:hotcold-space-correctness} 对比多层级内存结构、纯数据库结构和冷热分层结构。冷热分层在近期热点查询和近期删除负载下延迟较低，代价是空间占用高于纯内存紧凑结构。

本文实现了 k-kick 哈希表的简化原型，记录了与严格理论模型的差异，并给出测量数据与实验参数。冷热分层部分冷数据改为 SQLite3 记录后，近期访问更快，但失去了多层级 cubby 的紧凑性。

\section{不足与局限}

本文实验还有四方面不足：

(1) 实现的系统原型尚未与 IcebergHT\cite{pandey2023iceberght}、PaCHash\cite{kurpicz2023pachash} 等同类系统做同机对比，因此只能说明自身参数变化下的行为，难以直接评价相对性能；

(2) 未验证 k-kick 的树高，区间大小的取值等因素是否会影响系统性能效果；

(3) 第四章采用的合成负载较理想，规模覆盖也有限；

(4) 第五章更关注写后读和近期删除，对随机冷查询、数据库页面缓冲充分以及更复杂持久化策略的覆盖还不够。

\section{展望}

后续工作可在现有实验基础上进一步补充对照实验和运行指标。本文目前主要比较了基础多层级结构、纯 SQLite 结构与冷热分层结构的表现，后续可在相同硬件环境、相同数据规模和相同负载设置下，引入 PaCHash 或 IcebergHT 等紧凑哈希结构作为同机 baseline，从而更准确地判断本文方案在空间开销、查询延迟和更新代价上的相对位置\cite{kurpicz2023pachash,lehmann2026mphsurvey,pibiri2021pthash}。同时，现有实验主要关注单次操作延迟，对整体运行状态的刻画仍不充分，后续可继续记录进程级内存占用、SQLite 运行时文件大小、批量写入吞吐量以及混合负载下的查询和更新吞吐表现\cite{cooper2010ycsb}。在此基础上，还可以针对插入和删除热路径继续优化，并比较不同数据规模下 facility 规模、k-kick 深度和区间大小对性能的影响，以分析冷热分层结构在不同负载条件下是否存在更合适的参数组合。
